\rok{Новембар 2024.}{C}

Први практични тест из Алгоритама и структура података

\zadatak{1}{Quick Sort}
Написати \textit{quickSort} методу \textit{Quick Sort} алгоритма за сортирање целобројног низа дужине \textit{n}.

\textbf{Додатне напомене за решавање задатка:}
\begin{itemize}
    \item Улаз је несортиран низ целих бројева.
    \item Излаз је сортирани низ целих бројева.
    \item У шаблону се налази класа \texttt{RandomArrayGenerator} коју треба искористити за генерисање низа случајних бројева.
    \item У оквиру класе \texttt{QuickSort} написати имплементацију за методу:
    \begin{Verbatim}[fontsize=\footnotesize, numbers=left, xleftmargin=10mm]
private static void quickSort(int[] arr, int lower, int upper)
    \end{Verbatim}
    \item Обезбедити код у \texttt{Main} методи за тестирање алгоритма.
\end{itemize}



\zadatak{2}{Максималан збир близанаца}
У повезаној листи са \textbf{парним бројем чворова} величине $n$, дефинишемо \textbf{близаначки пар} чворова на следећи начин:

За \textbf{i-ти чвор} (почевши од индекса 0) постоји одговарајући \textbf{близанац} на позицији $(n-1-i)$, ако важи $0 <= i <= (n/2)-1$.

На пример:
\begin{itemize}
    \item Када је $n=4$, чвор са индексом 0 има близанца у чвору са индексом 3, док чвор са индексом 1 има близанца у чвору са индексом 2. Ова два пара су једини близаначки парови за листу величине 4.
\end{itemize}

\textbf{Збир близанаца} дефинише се као збир вредности чвора и његовог близанца. Задатак је да пронађете \textbf{највећи збир близаначких парова} у листи. Решење треба да пронађе максимални збир близанаца у $O(n)$ времену.

\textbf{Додатни захтеви за решавање задатка:}
\begin{itemize}
    \item Алгоритам писати у оквиру класе \texttt{LinkedList}.
    \item Захтеван потпис методе:
    \begin{Verbatim}[fontsize=\footnotesize, numbers=left, xleftmargin=10mm]
public int findMaxTwinSum()
    \end{Verbatim}
\end{itemize}

\textbf{Примери:}
\begin{itemize}
    \item \textbf{Пример 1:} \\
          \textbf{Улаз:} $5 \leftrightarrow 4 \leftrightarrow 2 \leftrightarrow 1$ \\
          \textbf{Излаз:} 6
    \item \textbf{Пример 2:} \\
          \textbf{Улаз:} $4 \leftrightarrow 2 \leftrightarrow 2 \leftrightarrow 3$ \\
          \textbf{Излаз:} 7
\end{itemize}



\zadatak{3}{Провера два реда}
За два реда \textit{queue1} и \textit{queue2} треба проверити да ли можемо добити \textit{queue2} из \textit{queue1} користећи само операције дефинисане у оквиру класе \texttt{Queue}. Елементи који се налазе у другом реду морају пратити исти редослед као у првом реду, али неки елементи из првог реда могу бити прескочени.

\textbf{Додатни захтеви за решавање задатка:}
\begin{itemize}
    \item Задатак решавати помоћу структуре реда (класа \texttt{Queue}) која је дата у оквиру шаблона задатка.
    \item Захтеван потпис методе:
    \begin{Verbatim}[fontsize=\footnotesize, numbers=left, xleftmargin=10mm]
public bool CanFormQueue(Queue queue2)
    \end{Verbatim}
\end{itemize}

\textbf{Примери:}
\begin{itemize}
    \item Red1: [1, 2, 3, 4, 5] \\
          Red2: [2, 3, 4] \\
          \textbf{Испис:} true
    \item Улазни низ: [1, 2, 3, 4, 5] \\
          Излазни низ: [1, 4, 2] \\
          \textbf{Испис:} false
\end{itemize}

\sablon{AISP2024_Test1_GrupaC}{2024/Novembar/GrupaC/AISP2024_Test1_GrupaC.linq}
